% SGA 7, tomo I — traducción española: portada, página de créditos e introducción.
% UTF-8 sin BOM; finales de línea LF. Fragmento; no se afirma una compilación autónoma.
% La transcripción de las páginas preliminares del impreso de 1972 (cubierta, guarda editorial, portadas,
% página de copyright, página de Internet Archive) se conserva en ORIGINAL_FRONT_PAGES_ES.tex, fuera del lector.
% Los avisos de laguna, obstrucción y procedencia son editoriales, no texto histórico.

\providecommand{\SGASevenSourceGapES}[1]{\textit{[laguna visual: #1]}}
\providecommand{\SGASevenSourceObstructionES}[1]{\textit{[oculto por #1]}}

% Portada de esta edición.
\thispagestyle{empty}
\begin{center}
\vspace*{2.5cm}
{\large SEMINARIO DE GEOMETRÍA ALGEBRAICA DE BOIS-MARIE\par}
\medskip
{\large 1967--1969\par}
\vspace{2.2cm}
{\Huge SGA 7 I\par}
\vspace{1.2cm}
{\LARGE Grupos de monodromía\par}
\smallskip
{\LARGE en geometría algebraica\par}
\vspace{2.2cm}
{\large Dirigido por A.~Grothendieck\par}
\medskip
{con la colaboración de M.~Raynaud y D.~S.~Rim\par}
\vspace{3.5cm}
{\large Traducción española\par}
\medskip
{\sisi{Edición diplomática}{Edición corregida}\par}
\end{center}
\clearpage

% Página de créditos (reverso de la portada).
\thispagestyle{empty}
\null\vfill
\begingroup\small
\noindent\textbf{Edición original.} \textit{Séminaire de Géométrie Algébrique du Bois-Marie 1967--1969 (SGA~7~I).
Groupes de monodromie en géométrie algébrique}, dirigé par A.~Grothendieck avec la collaboration de M.~Raynaud et
D.~S.~Rim. Lecture Notes in Mathematics~288, Springer-Verlag, Berlin, Heidelberg, Nueva York, 1972.
ISBN 3-540-05987-3. \copyright{}~Springer-Verlag Berlin Heidelberg 1972.\par
\medskip
\noindent\textbf{Esta edición.} Traducción al español producida por IA a partir del texto francés del impreso de
1972 (escaneo: Internet Archive, \url{https://archive.org/details/groupesdemonodro0288unse}). Las correcciones y la composición tipográfica también
están producidas por IA. \sisi{Esta edición diplomática reproduce el texto impreso con sus erratas; la edición
corregida las enmienda.}{Esta edición corregida enmienda las erratas del impreso; la edición diplomática reproduce
el texto tal como se imprimió.}\par
\medskip
\noindent Los derechos sobre la obra original pertenecen a sus autores o a los representantes o herederos de estos.
En la medida en que tengamos derechos sobre la transcripción, la traducción, las correcciones y la composición
tipográfica producidas por IA, estos se dedican al dominio público bajo CC0~1.0. Esta edición se publica con fines
educativos y sin ánimo de lucro, y no reclamamos ningún derecho que no poseamos. Si usted es titular de derechos sobre
esta obra y se opone a su publicación aquí, abra una incidencia en GitHub
(\url{https://github.com/KokunoYumeto/sga-es/issues}) y la retiraremos.\par
\endgroup
\clearpage

\pdfbookmark[0]{Introducción}{sga7.introduccion}

% SGA7I-FRONT-ES-P006 — página PDF rectora 7 — Introducción, V impreso inferido.
\begin{center}
INTRODUCCIÓN
\end{center}

Sean $S$ un esquema y $f:X\longrightarrow S$ un morfismo de esquemas. Si $f$ es
propio y liso, y el número primo $\ell$ es invertible sobre $S$, los grupos de
cohomología $\ell$-ádica $H^i(X_{\bar s},\mathbb Z_\ell)$ de las fibras geométricas de $f$ forman un sistema
local $\ell$-ádico sobre $S$. Si solo se supone que $f$ es propio, estos grupos son las fibras
de un haz $\ell$-ádico $R^i f_*\mathbb Z_\ell$ sobre $S$. La teoría de los ciclos evanescentes relaciona
la ramificación de este haz sobre $S$ con las singularidades de $f$.

Solo consideraremos el caso en que $S$ es un rasgo henseliano (= espectro de un
anillo de valoración discreta henseliano). En la práctica, este es el caso esencial. Remito
a (I 2.1) para una descripción heurística de la teoría. Sean $S=\operatorname{Spec}(V)$, $\bar\eta$
una clausura algebraica del cuerpo de fracciones $k(\eta)$ de $V$, y $\bar s$ la clausura algebraica
correspondiente del cuerpo residual $k(s)$ ($k(\bar s)$ es el cuerpo residual de la normalización
$\bar V$ de $V$ en $k(\bar\eta)$). En el caso particular en que $X$ es propio y plano sobre
$S$, de dimensión relativa $n$, y en que $f$ solo presenta un punto de no lisitud
$x\in X_s$, se definen $\operatorname{Gal}(\bar\eta,\eta)$-módulos $\phi^i$, de naturaleza puramente local en un entorno
de $x$, nulos para $i\notin[0,n]$ (I 4.2), y se construye una sucesión exacta larga
de $\operatorname{Gal}(\bar\eta,\eta)$-módulos

\[
\cdots \longrightarrow H^i(X_{\bar s},\mathbb Z_\ell)
\xrightarrow{\mathrm{sp}} H^i(X_{\bar\eta},\mathbb Z_\ell)
\longrightarrow \phi^i
\longrightarrow H^{i+1}(X_{\bar s},\mathbb Z_\ell)
\longrightarrow \cdots
\]

(el $\operatorname{Gal}(\bar s,s)$-módulo $H^i(X_{\bar s},\mathbb Z_\ell)$ se considera como un
$\operatorname{Gal}(\bar\eta,\eta)$-módulo con acción trivial de la inercia; $\mathrm{sp}$ es la flecha de especialización).

Se dan asimismo criterios, válidos por ahora solo en característica
$0$, para que los $\phi^i$, con $i$ pequeño, sean nulos (por ejemplo, si además $X$ es
localmente una intersección completa, $\phi^i=0$ para $i\ne n$ (I 4.6)).
\clearpage

% SGA7I-FRONT-ES-P007 — página PDF rectora 8 — Introducción, VI visible.
El teorema de monodromía afirma que la acción del subgrupo de inercia $I$
de $\operatorname{Gal}(\bar\eta,\eta)$ es cuasiunipotente: para $T\in I$, existen enteros $N>0$,
$M>0$ tales que el endomorfismo $(T^M-I)^N$ de $H^i(X_{\bar\eta},\mathbb Z_\ell)$ sea nulo. Se dan dos
demostraciones de este teorema. La primera (I.1.2), de naturaleza aritmética, se aplica
siempre que los grupos de cohomología considerados tengan propiedades de finitud razonables.
El punto clave es que, cuando $k(s)$ es de tipo finito sobre su subcuerpo
primo, la acción de $I$ es cuasiunipotente para toda representación $\ell$-ádica de
$\operatorname{Gal}(\bar\eta,\eta)$. La segunda demostración, más geométrica, requiere la pureza y la resolución
de singularidades: por ahora solo es válida en característica $0$
o para un $H^1$. Aporta información valiosa sobre el exponente de nilpotencia
$N$ ($N\leq i+1$ para un $H^i$). Como se verá más adelante, se presta bien,
sobre $\mathbb C$, a la comparación con la teoría trascendente.

Estos resultados, aplicados al $H^1$ de las variedades abelianas, es decir, a su módulo
de Tate, permiten estudiar su reducción módulo $p$. Se dan así dos
demostraciones del teorema de reducción estable de las variedades abelianas, según el cual,
tras una ramificación, la fibra especial conexa del modelo de Néron es una extensión de una
variedad abeliana por un toro. La primera (I. 6) retoma el método de la demostración
aritmética del teorema de monodromía. La segunda (IX 3.6) se apoya en un
análisis mucho más fino del modelo de Néron y de sus propiedades de polarización.

Las exposiciones I a V de Grothendieck no fueron redactadas. Se resumieron
en una exposición I. Los resultados allí enunciados se demuestran de manera sucinta, pero
esencialmente completa.

La exposición II aplica el método de los lápices de Lefschetz y (I 5.3) al
estudio del grupo fundamental. Para $S$ una superficie sobre un cuerpo algebraicamente cerrado
$k$, de exponente característico $p$, se demuestra que el grupo profinito $\pi_1^{(p)}(S,s)$, la completación
fuera de $p$ del grupo fundamental, es de pro-$(p)$-presentación finita.

Como se explicó anteriormente, las exposiciones III a V no existen.
\clearpage

% SGA7I-FRONT-ES-P008 — página PDF rectora 9 — Introducción, VII visible.
La exposición VI contiene, con algunos complementos, la teoría de deformaciones
de Schlessinger. Se procura tener en cuenta los automorfismos infinitesimales
de los objetos que se clasifican, y no pasar demasiado bruscamente al funtor de las clases
de isomorfía de objetos. Se da asimismo una nueva construcción de los
$\underline{\operatorname{Ext}}^{i}(L_{X/S},-)$ ($L_{X/S}$, complejo cotangente relativo de $X/S$). Esta exposición se utiliza en el
resto del seminario sobre todo a través del estudio que hace de las deformaciones de las singularidades
cuadráticas ordinarias.

Las exposiciones VII y VIII están dedicadas a la teoría de las biextensiones. Esta
teoría desempeña un papel esencial en la exposición IX para expresar qué ocurre con una
polarización al pasar de una variedad abeliana a uno de sus modelos de Néron.
Esta exposición IX contiene la demostración del teorema de reducción estable de las variedades
abelianas y diversas aplicaciones.

La continuación de este seminario, SGA 7 II, por P. Deligne y N. Katz, aparecerá
posteriormente.

\begin{flushright}
Bures sur Yvette, mayo de 1972,\\
P. DELIGNE
\end{flushright}
\clearpage

% SGA7I-FRONT-ES-P009 — página PDF rectora 10 — verso en blanco.
% No hay contenido impreso intencional. Solo se transparenta débilmente el reverso.
\null
\clearpage

% SGA7I-FRONT-ES-P010 — página PDF rectora 11 — Índice.
\begin{center}
Índice
\end{center}

\noindent
Exposición I\par
\hspace*{2em}Resumen de las primeras exposiciones de A. Grothendieck\par
\hspace*{2em}redactado por P. Deligne \dotfill 1\par

\medskip
\noindent
Exposición II\par
\hspace*{2em}Propiedades de finitud del grupo fundamental\par
\hspace*{2em}por Michèle Raynaud \dotfill 25\par

\medskip
\noindent
Exposición VI\par
% Título histórico impreso: Formal deformation theory.
\hspace*{2em}Teoría formal de deformaciones\par
\hspace*{2em}por D.S. Rim \dotfill 32\par

\medskip
\noindent
Exposición VII\par
\hspace*{2em}Biextensión de haces de grupos\par
\hspace*{2em}por A. Grothendieck \dotfill 133\par

\medskip
\noindent
Exposición VIII\par
\hspace*{2em}Complementos sobre las biextensiones. Propiedades generales\par
\hspace*{2em}de las biextensiones de esquemas en grupos\par
\hspace*{2em}por A. Grothendieck \dotfill 218\par

\medskip
\noindent
Exposición IX\par
\hspace*{2em}Modelos de Néron y monodromía\par
\hspace*{2em}por A. Grothendieck \dotfill 313\par
\clearpage

% SGA7I-FRONT-ES-P011 — página PDF rectora 12 — verso en blanco.
% No hay contenido impreso intencional. Solo se transparenta débilmente el reverso.
\null
\clearpage
